% !PRINTED 105
\weberpage{105}
%<<CONT

\begin{equation}
 |f(x)|>C.
\tag{3}
\end{equation}
To this is attached the theorem:

\medskip
\noindent V. \gesperrt{$f(x)$ is a continuous function of $x$.}
\medskip

By this the following is meant. If the absolute value of $x$ is smaller than an arbitrary finite magnitude $R$, then, however small the positive magnitude $\omega$ is assumed, one can find a positive magnitude $\varepsilon$ such that
\begin{equation}
 |f(x+h)-f(x)|<\omega
\tag{4}
\end{equation}
as long as the absolute value $\varrho$ of $h$ remains smaller than $\varepsilon$, and indeed in such a way that $\varepsilon$ depends only on the choice of $\omega$, not on $x$.

By § 13, (2),
\begin{equation}
\begin{split}
 f(x+h)-f(x)&=h f'(x)+\frac{h^2}{1\cdot2}f''(x)+\cdots+\frac{h^n}{\Pi(n)}f^{(n)}(x)\\
 &=h\left\{f'(x)+\frac{h}{2}f''(x)+\cdots+a_0h^{n-1}\right\}.
\end{split}
\tag{5}
\end{equation}
If now we assume the absolute value of $h$ smaller than $1$, then the absolute value of the magnitude in parentheses
\begin{equation}
 f'(x)+\frac{h}{2}f''(x)+\cdots+a_0h^{n-1}
\tag{6}
\end{equation}
is smaller than a certain non-zero number $k$ which one obtains if in (6) one replaces $h$ by $1$, the coefficients $a_0,a_1,\ldots,a_{n-1}$ by their absolute values and $x$ by $R$, because thereby every individual term of the sum (6) is replaced by its absolute value or by a larger number. If one chooses $\varepsilon$ so small that
\[
 \varepsilon k<\omega
\]
then by (5) the requirement (4) is fulfilled as long as
\[
 \varrho<\varepsilon
\]
remains.

We can give the theorem of the continuity of the function $f(x)$ a somewhat different expression, which is important for what follows. By the theorem that the absolute value of a sum of two terms lies, in magnitude, between the sum and the difference of the absolute values of the terms (cf. Introduction, p. 19), we obtain from (5), if we put
%CONT>>
